\section{Reference scales and rescaled probes}
\label{subsec:interacting-stop-probe}

We use the reference torus of the tuned flow and prove the
uniform probe estimates from \cite[\S21.1]{QED3Master}.
Fix $\bar m\in(0,1]$, put
$h=L^{\lfloor\log_L(1/\bar m)\rfloor}$, and assume
$0<e\sqrt h\le e_1$ and $N\ge N_{\rm ref}(e)$.
By \cite[(47)]{CoverTuning} and
\cite[(53)]{CoverTuning},
\begin{equation}\label{eq:interacting-physical-scales}
 \begin{aligned}
 \eta=L^{-N},\qquad K=N+\log_Lh,\qquad
 &a_j=L^{j-N},\\
 e_j=e\sqrt{a_j},\qquad
 &\eta_j=L^{-j}=\eta/a_j,\qquad \ell_j=R^*/a_j .
 \end{aligned}
\end{equation}
Thus a point with level-$j$ coordinates $c$
has physical coordinates $a_jc$, and
$T_j=\eta_j\Z^3/\ell_j\Z^3$.
At $j=0$ the charge is $q=e\sqrt\eta$.
The symbols $e_0,e_1$ without a running
index retain their meanings as charge ceilings.
Set $v_j=-\log e_j$, $p_{0,j}=v_j^{p_0}$ and $p_j=v_j^p$;
$r_j$ is the least nonnegative integer power of $L$ at least
$v_j^{r_{\rm loc}}$.
For a bond cochain $a$ on the fine unit torus,
its rescaling is
\begin{equation}\label{eq:interacting-connection-rescaling}
                    a^{[j]}(b)=L^{j/2}a(L^jb).
\end{equation}
The physical scale $a_j$ is a number, while
$a$ denotes a connection.
On the reference torus,
\cite[(54)]{CoverTuning} and
\cite[(50)]{CoverTuning} give
\[
 k_*\in\{K-1,K-2\},\qquad
 a_*\in\{h/L,h/L^2\},\qquad
 \ell_*=Mr_*,\qquad r_*>31L+16.
\]
These stop-scale quantities are independent
of $N\ge N_{\rm ref}(e)$.
Write $r_{\rm loc}$ for the localization
power in \cite[\S6.1]{CoverTuning}.

Let $\beta\in C_c^\infty(\R^3;\Lambda^2)$ be real
and nonzero. Use the cube
$Q_\beta=z_\beta+[-D_\beta^{\rm buf}/2,D_\beta^{\rm buf}/2]^3$
of \Cref{sec:reference-cells},
where $D_\beta^{\rm buf}=D_\beta+4$.
In this section a placement $v\in\eta\Z^3$
is admitted if, in physical coordinates,
\begin{equation}\label{eq:interacting-admitted-placement}
              v+Q_\beta\subset[\ell_c/8,7\ell_c/8]^3,
              \qquad \ell_c=R^*/(2L).
\end{equation}
A centered placement has
$|v+z_\beta-(\ell_c/2)(1,1,1)|_\infty\le\eta/2$.
It is admitted when
$D_\beta^{\rm buf}/2+\eta/2\le3\ell_c/8$.
This follows by adding the displacement
of the cube center to its half-side.

Form $\widehat{\mathfrak a}$ with $\tau_v\beta$
as in \eqref{eq:cells-midpoint-source}, and
periodize it on $C(n_*)$, $n_*=R^*/\eta$.
For $c\in\mathcal C_*=\{0,\ldots,2L-1\}^3$,
write $\widehat{\mathfrak a}_c$ for its
cell copy; thus $\widehat{\mathfrak a}_c=\widehat a_c$ in the
introduction. In physical coordinates the
wall map carrying the first cell to cell $c$
acts in coordinate $\mu$ by translation through
$2i\ell_c$ if $c_\mu=2i$, and by reflection
$y_\mu\mapsto(2i+2)\ell_c-y_\mu$ if $c_\mu=2i+1$.
Define the marked connections by
\begin{equation}\label{eq:interacting-marked-connections}
 a(t)=e^{-1}\sum_ct_c\widehat{\mathfrak a}_c,\qquad
 \widehat c_c=e^{-1}\widehat{\mathfrak a}_c^{[k_*]},\qquad
 \mathscr C(t)=a(t)^{[k_*]}=\sum_ct_c\widehat c_c .
\end{equation}
For a level-$j$ bond function $A$, $\partial A$
uses forward differences divided by $\eta_j$,
and
\[
 |\delta_\alpha\partial A|_\infty
  =\sup_{\mu,\nu}\sup_{0<d(x,x')\le1}
     \frac{|\partial_\nu A_\mu(x')-\partial_\nu A_\mu(x)|}
          {d(x,x')^\alpha}.
\]
The external norm of a fine connection is
$|a|_{\rm ext}=|a^{[k_*]}|_\infty+
|\partial a^{[k_*]}|_\infty+
|\delta_\alpha\partial a^{[k_*]}|_\infty$.

We define the uniform probe constants over both
possible stop scales. For $\eta'=L^{-N'}$ put
\[
 g_\nu^{\eta'}[\beta']
      =-\sum_{\mu\ne\nu}\delta_{\eta',\mu}J_{\eta'}\beta'_{\mu\nu}.
\]
This is the smooth vector field
$J_{\eta'}a'_{\eta'}$ in the midpoint formula.
Let $\mathcal I_\beta$ consist of triples
$(N',a_{\rm s},r')$ with
$N'\ge0$, $a_{\rm s}\in\{h/L,h/L^2\}$,
$N'+\log_La_{\rm s}\ge1$, and
$r'\in\{\pm1\}^3$.
$r'$ is identified with its diagonal orthogonal
matrix, and $U_{r'}\beta(x)=(\wedge^2r')\beta(r'^{-1}x)$.
At such a triple set $\eta''=\eta'/a_{\rm s}$,
$g^0=g^{\eta'}[U_{r'}\beta]$, and
\[
 G_\nu(c)=\sqrt{a_{\rm s}}\,
                  g_\nu^0(a_{\rm s}c+\eta'e_\nu/2)
                       \quad(c\in\eta''\Z^3),\qquad
 W_\nu(x)=\sqrt{\eta'}\,g_\nu^0(\eta'x+\eta'e_\nu/2)
                       \quad(x\in\Z^3).
\]
The three norms of $G$ are taken at mesh $\eta''$.
Set
\begin{equation}\label{eq:interacting-uniform-probe-constants}
 C_\beta=\sup_{\mathcal I_\beta}
              (|G|_\infty+|\partial G|_\infty+
                                    |\delta_\alpha\partial G|_\infty),
 \qquad
 \mathbf D=\sup_{\mathcal I_\beta}\|dW\|_{\ell^2(\text{plaquettes})}^2.
\end{equation}
These suprema depend on $\bar m,L,\alpha,\beta$,
and contain no charge, torus, or placement.
They have the finite bounds
\begin{equation}\label{eq:interacting-uniform-probe-finite}
 C_\beta\le2\sqrt h(1+h+3h^2)|\beta|_{C^3},
 \qquad
 \mathbf D\le48|\beta|_{C^2}^2(D_\beta^{\rm buf}+3)^3 .
\end{equation}

Indeed the centered difference is the average
of the corresponding derivative along its
segment, and $J_{\eta'}$ has norm one on
each derivative sup norm. Each component
of $g^0$ has two summands, hence
$|D^ig_\nu^0|_\infty\le2|\beta|_{C^{i+1}}$ for
$i=0,1,2$. This gives
$|G|_\infty\le2a_{\rm s}^{1/2}|\beta|_{C^1}$.
A first lattice difference of $G$ is
$a_{\rm s}^{3/2}$ times an average of
$\partial_\mu g_\nu^0$, so
$|\partial G|_\infty\le2a_{\rm s}^{3/2}|\beta|_{C^2}$.
The difference of these averages at $c,c'$
is bounded by
$6a_{\rm s}|c-c'|\,|\beta|_{C^3}$.
Since $|c-c'|\le1$ and $\alpha<1$,
the third norm is at most
$6a_{\rm s}^{5/2}|\beta|_{C^3}$.
Sum these bounds and use $a_{\rm s}\le h$.

For the energy, each curl component of $W$
is a difference of two first differences,
and has modulus at most
$4(\eta')^{3/2}|\beta|_{C^2}$.
A supporting plaquette has a supporting
bond, whose endpoints in physical coordinates
lie in $r'Q_\beta$.
Its basepoint is therefore in that cube
enlarged by $\eta'$ in each negative
coordinate direction. There are at most
$(D_\beta^{\rm buf}/\eta'+2)^3$ basepoints
and three plaquette directions. Squaring
the pointwise curl bound gives
$48|\beta|_{C^2}^2(D_\beta^{\rm buf}+2\eta')^3$,
which implies \eqref{eq:interacting-uniform-probe-finite}.

At the reference stop, $(N,a_*,1)\in\mathcal I_\beta$.
The midpoint formula and
\eqref{eq:interacting-connection-rescaling} give
\[
 \widehat{\mathfrak a}^{[k_*]}_\nu(c)
       =\sum_{z\in\Z^3}G_\nu(c-v/a_*+\ell_*z).
\]
Supporting bond endpoints lie in
$a_*^{-1}(v+Q_\beta)$.
Distinct periodic copies of that cube have
distance at least
$\ell_*-3\ell_*/(8L)>\ell_*/2>1+2\eta_*$.
Each functional defining one of the three
norms reads only bonds with endpoint distances
at most $1+2\eta_*$. It therefore reads
at most one periodic copy. Taking its
supremum identifies each torus norm with
the corresponding norm of $G$.
Translations and coordinate reflections
preserve these norms, including the one-bond
shift that reverses an oriented component.
The same argument applies to the cell copies.

The image of $v+Q_\beta$ under each wall map
lies in
$\prod_\mu[c_\mu\ell_c+\ell_c/8,(c_\mu+1)\ell_c-\ell_c/8]$.
For distinct cells, a separating coordinate
on the periodic circle gives distance
at least $\ell_c/4$: any joining arc
contains the two margins of size $\ell_c/8$.
Thus supporting endpoints of distinct
$\widehat c_c$ have stop-coordinate distance
at least $\ell_*/(8L)$, which is larger
than $1+2\eta_*$. Again each external-norm
functional reads at most one copy, so
\begin{equation}\label{eq:interacting-external-probe-bound}
 |a(t)|_{\rm ext}\le C_\beta|t|_\infty/e,\qquad
                 v_\eta=\|d\widehat{\mathfrak a}\|^2\le\mathbf D.
\end{equation}
For the second bound, the covering curvature
$dW$ is supported in the cube of side
$D_\beta^{\rm buf}/\eta+2<n_*$
obtained by enlarging $\eta^{-1}(v+Q_\beta)$.
Its periodic translates do not overlap,
so periodization preserves its counting
norm. Distinct cell curvatures also have
disjoint plaquette supports: a common
plaquette would have vertices in both
endpoint sets at physical distance at
most $\eta\le1$, whereas $\ell_c/4>1$.
This gives \eqref{eq:cells-orthogonality}
for every admitted placement.

Let $P_c$ be the unit-block centers whose
blocks contain an endpoint of a supporting
bond of $\widehat c_c$, and let $P_c^+$
be their $(M+1)$-neighborhood.
For $y\in P_c$, $y'\in P_{c'}$ choose the
corresponding endpoints in $\Delta_y,\Delta_{y'}$.
Each endpoint is at distance less than $1/2$
from its center. Consequently
\begin{align}
 d_{\rm sep}:=\min_{c\ne c'}d(P_c,P_{c'})
    &\ge\ell_*/(8L)-1=Mr_*/(8L)-1,
                              \label{eq:interacting-probe-separation}\\
 |P_c^+|&\le p_\beta:=
                  (D_\beta^{\rm buf}/a_*+2M+4)^3 .
                              \label{eq:interacting-support-number}
\end{align}
For the count, the endpoints have a lift in
a cube of side $D_\beta^{\rm buf}/a_*$.
A center adds less than $1/2$ to each
coordinate bound, and its neighborhood adds
$M+1$. The resulting open interval has
length $D_\beta^{\rm buf}/a_*+2M+3$
and contains fewer than one plus that
many integers. Multiply the three counts.
If the lattice probe vanishes, all these
sets are empty and $d_{\rm sep}=\infty$.

